% docs/manual/developer/chapters/06-core-engine.tex
% 第 6 章：core 评估引擎

\chapter{core 评估引擎}

\modname{core} 是 \openbench{} 的"算"层：拿到对齐好的 sim/ref 数据，算出 metrics、scores、comparison 表、各种统计。本章按子模块组织，每个关键扩展点配一个 walkthrough。

\section{子系统鸟瞰}

\modname{core} 包含：

\begin{itemize}
  \item \modname{metrics}：25+ 统计指标（bias、RMSE、KGE、NSE 等）
  \item \modname{scores}：8 个归一化打分（nBiasScore、Overall\_Score 等）
  \item \modname{evaluation}：Evaluation\_grid + Evaluation\_stn 两个评估类
  \item \modname{evaluation\_engine}：ModularEvaluationEngine + Grid/Station 子类
  \item \modname{comparison}：跨 simulation 合并指标
  \item \modname{landcover\_groupby} / \modname{climatezone\_groupby}：分组统计
  \item \modname{statistics/}：19 个 \texttt{stat\_*.py} 模块（ANOVA、相关、KDE 等）+ \modname{Mod\_Statistics} dispatcher
\end{itemize}

\section{metrics.py：统计指标}

\file{metrics.py} 定义 \texttt{class metrics}，每个方法是一个指标：

\begin{minted}{python}
class metrics:
    def __init__(self): ...
    def bias(self, s, o): ...
    def percent_bias(self, s, o): ...
    def RMSE(self, s, o): ...
    def ubRMSE(self, s, o): ...
    def correlation(self, s, o): ...
    def NSE(self, s, o): ...
    def KGE(self, s, o): ...
    # ... 25+ 个
\end{minted}

\subsection{方法约定}

\begin{itemize}
  \item 入参：\texttt{s}（simulation, numpy array）+ \texttt{o}（observation, numpy array），形状一致
  \item 返回：标量 float（grid 评估时是空间平均后的 metric；逐像素 metric 在外面调）
  \item NaN 处理：metric 内部用 \texttt{np.nanmean} 等 nan-aware 函数
  \item shape mismatch：一律 raise（由 evaluator 在更高层捕获）
  \item 不打日志（保持纯函数）
\end{itemize}

\subsection{Walkthrough：添加新 metric}

假设要加 \texttt{Pearson\_R3}（Pearson 相关的三次方，纯演示）：

\subsubsection{Step 1：在 metrics.py 加方法}

\begin{minted}{python}
def Pearson_R3(self, s, o):
    """Cubed Pearson correlation; emphasizes high agreement."""
    self._validate_inputs(s, o)
    r = self.correlation(s, o)
    return r ** 3
\end{minted}

\subsubsection{Step 2：加测试}

\begin{minted}{python}
# tests/test_core/test_metrics.py
def test_pearson_r3():
    m = metrics()
    s = np.array([1, 2, 3])
    o = np.array([1, 2, 3])
    assert m.Pearson_R3(s, o) == pytest.approx(1.0)
\end{minted}

\subsubsection{Step 3：注册到默认列表（可选）}

如果想让 \cli{openbench run} 默认计算它，把名字加到 \file{adapter.py} 的 default metrics 列表（或让用户在 yaml 显式指定）。

\subsubsection{Step 4：文档}

加到用户卷第 7 章的 metric 表 + 开发卷附录 B。

\section{scores.py：归一化打分}

\file{scores.py} 的 \texttt{class scores} 提供 8 个归一化打分：

\begin{itemize}
  \item \texttt{nBiasScore}：归一化偏差打分
  \item \texttt{nRMSEScore}：归一化均方根误差打分
  \item \texttt{nPhaseScore}：相位（季节循环）打分
  \item \texttt{nIavScore}：年际变化打分
  \item \texttt{nSpatialScore}：空间分布打分
  \item \texttt{nSeasonalityScore}：季节性打分
  \item \texttt{Overall\_Score}：综合分（其他分的加权平均）
\end{itemize}

\subsection{方法约定}

与 metrics 类似，但所有方法返回 [0, 1] 区间值（1 = 完美）。

\subsection{Walkthrough：添加新 score}

类似 metric。重点是归一化：把无界的 metric 转换到 [0, 1]，常用映射 \texttt{exp(-x)} 或 \texttt{1/(1+x)}。

\section{evaluation.py / evaluation\_engine.py}

\subsection{两层结构}

\begin{itemize}
  \item \modname{evaluation}：用户级 API（\texttt{Evaluation\_grid} / \texttt{Evaluation\_stn}）
  \item \modname{evaluation\_engine}：实现细节（\texttt{ModularEvaluationEngine} 基类 + Grid/Station 子类 + 工厂）
\end{itemize}

\subsection{ModularEvaluationEngine}

设计思路：每个评估单元（一对 sim×ref）由一个 engine 实例承担，engine 内部做：

\begin{enumerate}
  \item 调 \modname{data.processing.DatasetProcessing} 加载 + 对齐
  \item 调 \modname{metrics.metrics} 算所有指标
  \item 调 \modname{scores.scores} 算所有打分
  \item 写 CSV
  \item 调 \modname{visualization} 出图
\end{enumerate}

子类 \texttt{GridEvaluationEngine} 与 \texttt{StationEvaluationEngine} 区别在于站点形态的特殊处理。

\subsection{工厂函数}

\begin{minted}{python}
from openbench.core import create_evaluation_engine

engine = create_evaluation_engine(
    data_type="grid",   # 或 "stn"
    info={...},         # legacy info dict from adapter
)
result = engine.evaluate()
\end{minted}

\section{comparison.py}

跨 simulation 聚合：所有 sim 的 metric / score CSV 合并成对比表，输入到 portrait/radar/ranking 出图模块。

\subsection{触发}

仅当 \yamlkey{comparison.enabled: true}。\modname{runner.local} 在 evaluation phase 完成后调它。

\subsection{产物}

\begin{itemize}
  \item \file{output/<run>/comparisons/portrait\_table.csv}：模型 × (变量, score) 矩阵
  \item \file{output/<run>/comparisons/ranking.csv}：按 Overall\_Score 排序
  \item 调 \modname{visualization} 出 portrait\_plot.png + radar\_plot.png + ranking\_table.png
\end{itemize}

\section{Groupby 模块}

\modname{landcover\_groupby} 与 \modname{climatezone\_groupby}：

\begin{itemize}
  \item 输入：grid 评估的 metric/score 数据 + 一个分类掩码（IGBP / PFT / Köppen）
  \item 操作：按类别分组，每组算 metric/score 平均
  \item 输出：每分类一行的 CSV + 分类直方图（LC\_based\_heat\_map）
\end{itemize}

\subsection{触发}

仅当 \yamlkey{IGBP\_groupby: true} / \yamlkey{PFT\_groupby: true} / \yamlkey{climate\_zone\_groupby: true}。

\subsection{依赖外部数据}

\begin{itemize}
  \item IGBP 类掩码：\openbench{} 自带（在 \file{data/} 下）
  \item Köppen 气候带掩码：自带
  \item PFT 掩码：取决于具体 model（CoLM / CLM 等各有版本）
\end{itemize}

\section{statistics/：19 个统计模块}

\file{core/statistics/} 下：

\begin{itemize}
  \item \modname{Mod\_Statistics}：dispatcher，按 \yamlkey{statistics.items} 分发到具体模块
  \item \modname{stat\_anova}、\modname{stat\_correlation}、\modname{stat\_covariance}
  \item \modname{stat\_hellinger\_distance}、\modname{stat\_kernel\_density\_estimate}
  \item \modname{stat\_mann\_kendall\_trend\_test}、\modname{stat\_three\_cornered\_hat}
  \item \modname{stat\_partial\_least\_squares\_regression}、\modname{stat\_functional\_response}
  \item \modname{stat\_resample}、\modname{stat\_rolling}、\modname{stat\_standard\_deviation}、\modname{stat\_variance}
  \item \modname{stat\_z\_score}、\modname{stat\_diff}、\modname{stat\_autocorrelation}
  \item \modname{stat\_False\_Discovery\_Rate}、\modname{stat\_Basic}
  \item \modname{base}：共享基类与工具
\end{itemize}

\subsection{命名约定}

\begin{itemize}
  \item 文件名：\texttt{stat\_<name>.py}（snake\_case）
  \item 主函数名：与文件同名，例如 \texttt{stat\_anova.run\_anova(...)}
  \item I/O：从 \texttt{info} dict 读路径，写 CSV 到 \file{output/<run>/statistics/<name>/}
  \item 出图：调对应的 \modname{visualization.Fig\_<Name>} 模块
\end{itemize}

\subsection{Walkthrough：添加新 statistic 模块}

假设要加"Theil-Sen 斜率检验"：

\subsubsection{Step 1：建文件 \file{stat\_theil\_sen.py}}

\begin{minted}{python}
"""Theil-Sen non-parametric trend test."""
from .base import BaseStatistic

class TheilSenTest(BaseStatistic):
    name = "Theil_Sen"

    def compute(self, info):
        # 加载数据 + 计算 + 返回 dict
        ...

def run_theil_sen(info):
    return TheilSenTest().run(info)
\end{minted}

\subsubsection{Step 2：注册到 dispatcher}

\file{Mod\_Statistics.py} 中加：

\begin{minted}{python}
from . import stat_theil_sen

DISPATCH = {
    # ...
    "Theil_Sen": stat_theil_sen.run_theil_sen,
}
\end{minted}

\subsubsection{Step 3：加 visualization}

\file{visualization/Fig\_Theil\_Sen.py}：

\begin{minted}{python}
def make_Theil_Sen(info, result):
    """Plot Theil-Sen results to figures/Theil_Sen/."""
    ...
\end{minted}

\subsubsection{Step 4：在 yaml 中启用}

\begin{minted}{yaml}
statistics:
  enabled: true
  items:
    - Theil_Sen
\end{minted}

\subsubsection{Step 5：测试}

\file{tests/test\_core/test\_statistics/test\_theil\_sen.py}。

\section{Evaluation 内的并行}

\file{processing.py} 的 DatasetProcessing 在两个层次用 \modname{joblib} 并行：

\begin{itemize}
  \item 站点级：每个站点独立处理（用于 stn 评估）
  \item 按年级：多年 NC 合并时按年并发（grid 评估）
\end{itemize}

变量级（task 级）并发尚未启用——这是 v3.0 的已知 backlog。

\section{添加新 metric / score / statistic 的总结对照}

\begin{longtable}{l p{3cm} p{6cm}}
  \toprule
  扩展类型 & 加在哪 & 步骤 \\
  \midrule
  \endhead
  metric & \file{core/metrics.py} 内的 \texttt{class metrics} & 加 method + 加测试 + 文档（5 步） \\
  score & \file{core/scores.py} 内的 \texttt{class scores} & 同上但返回 [0, 1] \\
  statistic & 新建 \file{core/statistics/stat\_X.py} & 加文件 + 注册到 \modname{Mod\_Statistics} + 加 visualization + 测试 + 文档（5 步） \\
  comparison metric & \file{core/comparison.py} & 改 portrait\_table 计算逻辑 \\
  groupby & 新建 \file{core/<X>\_groupby.py} & 较少见；通常已有 IGBP/PFT/CZ 够用 \\
  \bottomrule
\end{longtable}

\section{下一步}

下一章详解 \modname{visualization}：17 个 Fig 模块的命名约定 + 添加新 Fig 模块的 walkthrough。
